% 自動生成（tools/mkchecklist.py p4）。各章の「ふりかえり」の項目を集めたもの（第13回）
\begingroup\footnotesize
\def\bkcb{\tikz[baseline=-0.5ex]\draw[\bkTF, line width=0.7pt, rounded corners=0.5pt] (0,-0.3em) rectangle (0.8em,0.5em);}
\begin{longtblr}{colspec={Q[c,m,wd=11mm] X[l,m] Q[c,m,wd=10mm] Q[c,m,wd=9mm] Q[c,m,wd=9mm]}, width=\linewidth, rowhead=1, colsep=3pt,
  row{1}={bg=bkVol!14, font=\sffamily\bfseries\scriptsize}, hlines={0.4pt,bkLine}, rowsep=0.7pt}
章 & できるようになったこと & ページ & 1回め & 2回め \\
{\sffamily\bfseries\color{\bkTT}電-0} & 発電所から家のドライヤーまで、電気が通る道すじと、とちゅうで働く機械（発電機・変圧器・送電線）を説明できる。 & \pageref{ph:電-0:sec1} & \bkcb & \bkcb \\
 & 何を無視したかを書いてから、電流・電力・電線で失われる電力を、7つの段で求められる。 & \pageref{ph:電-0:ex3} & \bkcb & \bkcb \\
 & 身近な「なぜ？」の答えが、この巻のどの章にあるかを探せる。 & \pageref{ph:電-0:sec5} & \bkcb & \bkcb \\
{\sffamily\bfseries\color{\bkTT}電-1} & こすると物が電気を帯びるわけを、電子が移ることで説明できる。 & \pageref{ph:電-1:sec1} & \bkcb & \bkcb \\
 & クーロンの法則と重ね合わせで、いくつかの電荷のあいだの力と、電場を求められる。 & \pageref{ph:電-1:sec3} & \bkcb & \bkcb \\
 & 一様な電場の中の電荷の動きを、力学の放物運動と同じ考えで求められる。 & \pageref{ph:電-1:sec8} & \bkcb & \bkcb \\
 & 線・円環・平面に広がった電荷の電場を積分で求め、電気双極子の電場が距離の3乗に反比例して弱まることを示せる。 & \pageref{ph:電-1:sec7} & \bkcb & \bkcb \\
{\sffamily\bfseries\color{\bkTT}電-2} & ガウスの法則で、球・長い直線・広い平面の電荷の電場を求められる。 & \pageref{ph:電-2:sec4} & \bkcb & \bkcb \\
 & 電位を使って、電荷を運ぶ仕事や、加速された電子の速さを求められる。 & \pageref{ph:電-2:sec5} & \bkcb & \bkcb \\
 & 等電位面と電気力線の関係を読み、一様な電場で $V=Ed$ を使える。 & \pageref{ph:電-2:sec7} & \bkcb & \bkcb \\
 & ガウスの法則の微分形と $\vec E=-\mathrm{grad}\,\phi$ を、ベクトル解析で説明できる。 & \pageref{ph:電-2:sec8} & \bkcb & \bkcb \\
 & ラプラス方程式を表計算で解き、電場が集まる所を示せる。 & \pageref{ph:電-2:sec11} & \bkcb & \bkcb \\
{\sffamily\bfseries\color{\bkTT}電-3} & コンデンサーにためた電気で、豆電球や LED を光らせられる。 & \pageref{ph:電-3:sec1} & \bkcb & \bkcb \\
 & 導体の中の電場が0になるわけと、静電遮へいを説明できる。 & \pageref{ph:電-3:sec3} & \bkcb & \bkcb \\
 & 平行板の静電容量・つなぎ方・エネルギーを求められる。 & \pageref{ph:電-3:sec4} & \bkcb & \bkcb \\
 & 分極と電束密度で、同軸・球の静電容量を導ける。 & \pageref{ph:電-3:sec8} & \bkcb & \bkcb \\
 & 境界条件・仮想変位・鏡像法を使って、電場や力を求められる。 & \pageref{ph:電-3:sec9} & \bkcb & \bkcb \\
{\sffamily\bfseries\color{\bkTT}電-4} & オームの法則・直列と並列・電力 $P=VI$ で、回路を調べられる。 & \pageref{ph:電-4:sec3} & \bkcb & \bkcb \\
 & 抵抗率から抵抗を求め、電子の流れとジュール熱を説明できる。 & \pageref{ph:電-4:sec5} & \bkcb & \bkcb \\
 & 抵抗の温度変化、電池の内部抵抗、キルヒホッフの法則を使える。 & \pageref{ph:電-4:sec7} & \bkcb & \bkcb \\
 & $\vec j=\sigma\vec E$・連続の式から、オームの法則と電流則を導ける。 & \pageref{ph:電-4:sec6} & \bkcb & \bkcb \\
{\sffamily\bfseries\color{\bkTT}電-5} & 乾電池と豆電球の直列・並列で、電流と明るさがどう変わるかを説明できる。 & \pageref{ph:電-5:sec2} & \bkcb & \bkcb \\
 & 電流の分かれ方・ブリッジ・電流計と電圧計・コンデンサーをふくむ回路を、キルヒホッフの法則で計算できる。 & \pageref{ph:電-5:sec1} & \bkcb & \bkcb \\
 & 網目電流法・節点電位法で回路を行列の式にし、テブナンの定理や Python で解ける。 & \pageref{ph:電-5:sec3} & \bkcb & \bkcb \\
{\sffamily\bfseries\color{\bkTT}電-6} & 電流の周りに輪の形の磁場ができることを方位磁針で確かめ、その向きを右ねじで言える。 & \pageref{ph:電-6:sec2} & \bkcb & \bkcb \\
 & まっすぐな電流・円い電流・ソレノイドがつくる磁場と、磁場の中の電流・荷電粒子・コイルが受ける力を求められる。 & \pageref{ph:電-6:sec4} & \bkcb & \bkcb \\
 & ビオ・サバールの法則とアンペールの法則（積分形）から磁場を導き、$\mathrm{rot}\,\vec B=\mu_0\vec j$ と $\mathrm{div}\,\vec B=0$ を読める。 & \pageref{ph:電-6:sec6} & \bkcb & \bkcb \\
 & ホール効果・速度選別器・サイクロトロンのしくみを、ローレンツ力で説明できる。 & \pageref{ph:電-6:sec10} & \bkcb & \bkcb \\
{\sffamily\bfseries\color{\bkTT}電-7} & 磁石に近づけた鉄が磁石になるわけを、小さな磁石の向きがそろう図で説明できる。 & \pageref{ph:電-7:sec2} & \bkcb & \bkcb \\
 & 磁化 $M$・磁場の強さ $H$・比透磁率を使って、鉄心の中の磁場を求められる。 & \pageref{ph:電-7:sec3} & \bkcb & \bkcb \\
 & 磁気回路を電気回路と同じ形で解き、すき間のある電磁石の磁束を求められる。 & \pageref{ph:電-7:sec5} & \bkcb & \bkcb \\
 & 鉄と空気の境目で磁力線がどう曲がるかを、境界条件から求められる。 & \pageref{ph:電-7:sec7} & \bkcb & \bkcb \\
 & $BH$ 曲線の面積からヒステリシス損を、減磁曲線から永久磁石の動作点を求められる。 & \pageref{ph:電-7:sec8} & \bkcb & \bkcb \\
{\sffamily\bfseries\color{\bkTT}電-8} & 磁石をコイルに出し入れすると電流が流れ、動かす向きと速さで、電流の向きと大きさが変わることを、実験で確かめられる。 & \pageref{ph:電-8:sec1} & \bkcb & \bkcb \\
 & ファラデーの法則とレンツの法則で、動く棒や発電機に生まれる電圧（誘導起電力）の大きさと向きを求められる。 & \pageref{ph:電-8:sec3} & \bkcb & \bkcb \\
 & 自己インダクタンス・相互インダクタンスを求め、コイルにたまるエネルギーを計算できる。 & \pageref{ph:電-8:sec6} & \bkcb & \bkcb \\
 & ファラデーの法則の微分形を読み、磁場のエネルギーから電磁石が鉄を引く力を求められる。 & \pageref{ph:電-8:sec9} & \bkcb & \bkcb \\
 & 渦電流と表皮効果で、IH調理器や、うすい板を重ねた鉄心のしくみを説明できる。 & \pageref{ph:電-8:sec12} & \bkcb & \bkcb \\
{\sffamily\bfseries\color{\bkTT}電-9} & スイッチを入れた瞬間と、十分に時間がたったあとの回路を、コンデンサーとコイルの性質から説明できる。 & \pageref{ph:電-9:sec1} & \bkcb & \bkcb \\
 & 電圧則と部品のきまりから RC・RL 回路の式を立て、時定数と、エネルギーの行き先を求められる。 & \pageref{ph:電-9:sec3} & \bkcb & \bkcb \\
 & LC 回路の電気振動の周波数とエネルギーの行き来、RLC 回路の臨界抵抗を求められる。 & \pageref{ph:電-9:sec7} & \bkcb & \bkcb \\
 & ラプラス変換と数値計算で RLC 回路の過渡現象を求め、力学と電気の相似の表を、電気の法則から説明できる。 & \pageref{ph:電-9:sec9} & \bkcb & \bkcb \\
{\sffamily\bfseries\color{\bkTT}電-10} & 家の交流の周波数と実効値の意味を説明し、\qty{100}{V}の交流の山の電圧を求められる。 & \pageref{ph:電-10:sec2} & \bkcb & \bkcb \\
 & コイル・コンデンサーのリアクタンスと位相のずれ、RLC 直列回路のインピーダンスと共振周波数を求められる。 & \pageref{ph:電-10:sec4} & \bkcb & \bkcb \\
 & 電圧・電流を複素数（フェーザ）で表して交流回路を解き、Q 値・フィルター・交流ブリッジを計算できる。 & \pageref{ph:電-10:sec5} & \bkcb & \bkcb \\
{\sffamily\bfseries\color{\bkTT}電-11} & 交流の電力が時間とともにゆれること、家の電気と送電線の電気の違いを、図で説明できる。 & \pageref{ph:電-11:sec2} & \bkcb & \bkcb \\
 & 力率を使って、有効電力・無効電力・皮相電力を計算できる。 & \pageref{ph:電-11:sec4} & \bkcb & \bkcb \\
 & コンデンサーで力率をよくすると、電線の電流と損失が減るわけを説明できる。 & \pageref{ph:電-11:sec5} & \bkcb & \bkcb \\
 & 三相交流の Y 結線・$\Delta$ 結線の電圧と電流、三相の電力を計算できる。 & \pageref{ph:電-11:sec7} & \bkcb & \bkcb \\
 & 三相の電力が時間によらず一定になること、回転磁界ができることを、式で示せる。 & \pageref{ph:電-11:sec9} & \bkcb & \bkcb \\
{\sffamily\bfseries\color{\bkTT}電-12} & 磁石とコイルでモーターが回るわけを、電流が磁場から受ける力で説明できる。 & \pageref{ph:電-12:sec1} & \bkcb & \bkcb \\
 & 直流モーターのトルクと回る速さの関係を、逆起電力を使って説明できる。 & \pageref{ph:電-12:sec4} & \bkcb & \bkcb \\
 & 変圧器の電圧と電流の比を、巻数から求められる。 & \pageref{ph:電-12:sec5} & \bkcb & \bkcb \\
 & 変圧器の鉄心の磁束、損失と効率を計算できる。 & \pageref{ph:電-12:sec6} & \bkcb & \bkcb \\
 & 誘導モーターのすべりとトルク、同期モーターとインバーターのしくみを説明できる。 & \pageref{ph:電-12:sec7} & \bkcb & \bkcb \\
{\sffamily\bfseries\color{\bkTT}電-13} & ケーブルの上を電気の波が進み、先やつなぎ目ではね返ることを、速さ・往復の時間・波長で説明できる。 & \pageref{ph:電-13:sec1} & \bkcb & \bkcb \\
 & 特性インピーダンスの違いから、電圧の反射係数と、はね返る電力の割合を求められる。 & \pageref{ph:電-13:sec1} & \bkcb & \bkcb \\
 & 四端子定数の行列で箱をつなぎ、短い配電線の電圧降下と、変圧器の短絡電流を見積もれる。 & \pageref{ph:電-13:sec3} & \bkcb & \bkcb \\
 & 小さな箱のくさりから電信方程式を導き、波の速さ $1/\sqrt{L'C'}$ と特性インピーダンス $\sqrt{L'/C'}$ を求められる。 & \pageref{ph:電-13:sec6} & \bkcb & \bkcb \\
 & 損失のある線・長い送電線・4分の1波長の線を、$\cosh$・$\sinh$ と四端子定数で扱える。 & \pageref{ph:電-13:sec9} & \bkcb & \bkcb \\
{\sffamily\bfseries\color{\bkTT}電-14} & 充電しているコンデンサーの極板の間を流れる変位電流と、その周りの磁場を計算できる。 & \pageref{ph:電-14:sec3} & \bkcb & \bkcb \\
 & マクスウェル方程式を積分形と微分形で読み、電磁場のエネルギーの流れ（ポインティングベクトル）を計算できる。 & \pageref{ph:電-14:sec7} & \bkcb & \bkcb \\
 & 動く人から見ると電場と磁場が入れかわることを説明し、電磁場を数値で解く方法（FDTD）を表計算で動かせる。 & \pageref{ph:電-14:sec9} & \bkcb & \bkcb \\
{\sffamily\bfseries\color{\bkTT}電-15} & 電波・光・X 線などが、波長（周波数）の違う同じ電磁波であることを知り、$c=f\lambda$ で波長を計算できる。 & \pageref{ph:電-15:sec1} & \bkcb & \bkcb \\
 & 電磁波が電場と磁場の直角にゆれる横波で、$E=cB$ の関係をもち、エネルギーと運動量を運ぶことを計算で確かめられる。 & \pageref{ph:電-15:sec4} & \bkcb & \bkcb \\
 & ゆれる電荷が出す電磁波（ラーモアの式）、物の表面での反射率、金属の網が電波を通さないわけを、マクスウェル方程式から説明できる。 & \pageref{ph:電-15:sec8} & \bkcb & \bkcb \\
{\sffamily\bfseries\color{\bkTT}電-16} & 発電所から家までの電気の道すじと、東と西で周波数が違うわけ・東西をつなぐしくみを説明できる。 & \pageref{ph:電-16:sec1} & \bkcb & \bkcb \\
 & 整流・インバーター・直流送電を、スイッチの切りかえと「平均の値」で説明し、整流した電圧の平均を計算できる。 & \pageref{ph:電-16:ex2} & \bkcb & \bkcb \\
 & つくる電気と使う電気がずれると周波数が変わるわけを、回る発電機の運動エネルギーで説明し、周波数の下がる速さを見積もれる。 & \pageref{ph:電-16:sec2} & \bkcb & \bkcb \\
 & 調速機の働きを入れた周波数の変わり方を、プログラムで計算して読める。 & \pageref{ph:電-16:prog} & \bkcb & \bkcb \\
\end{longtblr}
\endgroup
